\newcommand{\figuraCapilarPlacas}{%
\begin{figure}[H]
\centering
\begin{tikzpicture}[>=Latex,font=\small,line cap=round,line join=round]
  \fill[blue!10] (0.7,0.55) rectangle (11.3,2.1);
  \draw[very thick,blue!65!black] (0.7,2.1)--(11.3,2.1);
  \draw[very thick] (4.65,0.55)--(4.65,5.35) (7.35,0.55)--(7.35,5.35);
  \fill[blue!18] (4.73,0.63)--(7.27,0.63)--(7.27,4.05)
    .. controls (6.65,3.65) and (5.35,3.65) .. (4.73,4.05)--cycle;
  \draw[very thick,blue!70!black] (4.73,4.05)
    .. controls (5.35,3.65) and (6.65,3.65) .. (7.27,4.05);
  \draw[<->,red!75!black,thick] (7.85,2.1)--(7.85,3.88)
    node[midway,right] {$h$};
  \draw[<->,thick] (4.65,5.65)--(7.35,5.65) node[midway,above] {$W$};
  \draw[->,red!75!black] (4.73,4.05) arc[start angle=90,end angle=30,radius=0.55];
  \node[red!75!black] at (5.18,4.50) {$\theta$};
  \draw[->,ultra thick,purple!75!black] (4.73,4.05)--++(120:0.95)
    node[above left] {$\alpha$};
  \draw[->,ultra thick,purple!75!black] (7.27,4.05)--++(60:0.95)
    node[above right] {$\alpha$};
  \draw[->,thick,dashed,red!75!black] (7.27,4.05)--(7.27,4.88)
    node[midway,left] {$\alpha\cos\theta$};
  \node[align=center] at (9.8,4.2)
    {dos líneas de contacto\\[1mm]sostienen la columna};
\end{tikzpicture}
\caption{Ascenso capilar entre placas paralelas separadas una distancia $W$.
Las componentes verticales de la tensión superficial en ambas paredes
equilibran el peso del líquido elevado.}
\label{fig:capilar-placas}
\end{figure}%
}